\section{Starting point}
I started this project after reading Exa's description of its vector database~\cite{exa}.
If we look at the work inside a selected cluster, one thing we can explore is whether the
query can help us choose documents before we score them. We can always calculate a score
for every document and return the best ones, but that gets expensive as the collection grows.
Skipping documents might help, so long as the work of choosing them costs less than the scores
we avoid, and we still find the results we wanted.

We'll compare three ways to search the same binary document embeddings. In method A, we open
selected clusters and scan their documents. In B, we open clusters and then split the documents
using bitplanes. In C, we open clusters and look up bit patterns, starting with the pattern
the query prefers. We'll call C \emph{weighted key lookup}; this is the concrete version of
starting at the ideal answer and ``walking backwards.''

If we fix the document count, RAM budget, and required recall, which method gives us the lowest
latency? Recall tells us how many of the correct reference results we recover, so we need to
keep it in the comparison as we change the search. And if we want to estimate latency, counting
calculations is only part of the job: we also need to look at the memory those calculations
read. Avoiding a document score can introduce work elsewhere.

The experiments give us three different outcomes to explain. On the tested real embeddings,
we did not find reliable evidence that either alternative was faster than a well-tuned scan. With constructed
queries whose important coordinates stayed in one prefix, key lookup had a small advantage
over a closely related directly routed scan. When the important coordinates changed with
each query, bitplane branching had a large advantage in the constructed case. So the useful
question is what makes those cases different, rather than which method wins everywhere.

Let's start with one query and six documents, follow it through each method, and then count
the work and storage. From there, we can look at recall, choose parameters, and compare the
measurements. At the end we'll use the calculations to consider a billion documents. Those
larger latency numbers are forecasts; the complete C++ indexes were tested through eight
million documents.

\subsection{A query example}
We use the binary document representation and floating-point query described in Exa's
post~\cite{exa}. In our examples, a stored bit of 1 means $+1$, and a bit of 0 means $-1$. We'll use five coordinates so we can follow
the calculation on the page. Let the query be
\[
q=[0.7,\ 0.5,\ -0.4,\ -0.6,\ 0.3].
\]
With that same dot-product score, signs \texttt{11001} make every contribution positive, giving
$0.7+0.5+0.4+0.6+0.3=2.5$. This is the pattern the query would prefer. Next we'll compare it
with the documents that actually exist.

\begin{center}
\begin{tabular}{@{}rcrr@{}}\toprule
Document ID & Stored signs & First-bit cluster & Full score\\\midrule
0 & \texttt{01001} & 0 & 1.1\\
1 & \texttt{00100} & 0 & $-1.3$\\
2 & \texttt{10100} & 1 & 0.1\\
3 & \texttt{11001} & 1 & 2.5\\
4 & \texttt{11000} & 1 & 1.9\\
5 & \texttt{00011} & 0 & $-1.1$\\\bottomrule
\end{tabular}
\end{center}
Here the correct top two are documents 3 and 4. If scores tie, we put the smaller document ID
first, and then use that same rule when checking every method. We can also run this example
against the C++ implementations with \path{examples/small_search.py}.

If we look at coordinate 1, a matching sign contributes $+0.7$, while a mismatch contributes
$-0.7$. That one change lowers the score by $1.4$. At coordinate 5, the loss is only $0.6$,
so the query magnitude tells us how expensive a wrong sign would be. We can use that observation
to guide the search. For a query with $d$ coordinates, write its ideal score as
$Q=\sum_{i=1}^{d}|q_i|$. Writing $b_i$ for the decoded sign ($+1$ or $-1$), if the mismatching
coordinates have total magnitude $p$, then
\begin{equation}
S(q,b)=Q-2p,
\qquad p=\sum_{i:\,b_i\ne\operatorname{sign}(q_i)}|q_i|.
\label{eq:score}
\end{equation}
For document 4, only coordinate 5 mismatches, giving $S=2.5-2(0.3)=1.9$. From here, ranking
by the largest score or the smallest mismatch penalty gives the same order. Ordinary Hamming
distance counts each mismatch equally; here we use the query magnitudes as weights.

\subsection{Shared pipeline}
Suppose the collection contains $N$ documents. Before any queries arrive, we embed the text,
choose a representation width $d$, pack the document signs, and build the index. We also assign
each document to one cluster. A cluster is a stored subset of the collection, and we can form
those subsets by learning from the vectors or by using a fixed sign prefix.

For shortening embeddings, we also refer to Exa's use of Matryoshka training~\cite{exa,mrl}.
We'll use that representation as given. One distinction matters for the proposed search:
a useful prefix is not the same thing as the largest query coordinates. Matryoshka doesn't
sort each query by magnitude, and it doesn't guarantee that its first 24 sign bits preserve
its nearest neighbors. Shortening the vector and taking its signs are two separate changes;
we hold the resulting representation fixed across the methods we compare.

Once a query arrives, we embed and normalize it, select clusters, and search their documents.
Opening one cluster is one \emph{probe}. With $C$ clusters and $P$ probes, balanced clusters
would give us about $NP/C$ selected documents. If the clusters aren't balanced, we need to
add their actual sizes instead. And if we want to search the whole collection, we can set
$C=P=1$.

\fig{shared}{The common query flow. The three local search methods share the same query and selected clusters when we isolate their work. Longer-vector reranking is optional; it is excluded from the primary binary-score comparison.}{.68}

A production pipeline can also filter documents or rerank candidates with longer vectors or
another model. If those stages are enabled, they belong in the complete request measurement.
For the primary index comparison here, we use cached query embeddings, the same binary score,
and no float reranker. Later we'll measure actual text queries so we can see the complete wait.

The baseline therefore follows Exa's public description of compact scoring inside selected
clusters~\cite{exa}. The measurements here are of our local implementation, not Exa's production
system. We also try two ways of forming clusters: learning groups from similar vectors,
known as IVF, and assigning documents by their first few sign bits. With a $k$-bit prefix,
there are up to $2^k$ possible patterns, not $k$ clusters. B and C change how we search the
selected documents; they do not introduce the shared embedding or quantization techniques.